\section{Related Work}
\label{sec:related}

\subsection{Static and Dynamic Windows Malware Detection}

Static malware analysis derives evidence without running a suspect object.
Learned detectors have used engineered binary attributes
\cite{saxe2015deep}, complete byte streams \cite{raff2018malware}, and documented
PE representations such as EMBER \cite{anderson2018ember}. SOREL-20M expanded
the scale of static PE research and supplied richer metadata
\cite{harang2020sorel}. Gradient-boosted trees remain effective on heterogeneous
tables containing histograms, strings, imports, sections, and header fields
\cite{ke2017lightgbm}. Their principal operational benefit is early scoring, but
static associations can be changed through packing, obfuscation, build-tool
variation, and distribution shift \cite{ucci2019survey,gibert2020rise}.

Dynamic analysis instead observes execution. API calls, file and registry
operations, process creation, and network activity can expose behavior absent
from a static vector \cite{sikorski2012practical,sgandurra2016automated}. The
evidence is conditional on the sandbox, inputs, run duration, and path reached by
the program. Delayed or environment-aware behavior can remain unobserved, while
longer collection increases resource cost. Work on ransomware illustrates the
value of behavior-centric analysis \cite{kharraz2015cutting,bacci2021ransomware},
but the present datasets contain generic malware labels. Those studies motivate
the analysis modes; they do not transfer ransomware semantics to our labels.

Hybrid systems commonly concatenate features, combine decisions, or fuse
scores. A concatenated model assumes that all modalities are available at once,
whereas sequential monitoring produces dynamic prefixes over time. Fixed score
averaging is transparent but cannot respond to sample-specific disagreement.
Learned stacking can capture a calibrated relation between channels, but it
requires a representative development cohort and can obscure how the weight
changes as evidence arrives. Meaningful evaluation must compare adaptive rules
with simple average, max-risk, fixed weighting, and a learned stacker using the
same paired samples and operating-point procedure.

\subsection{From Physical PID Control to Score Fusion}

PID control is a foundational feedback mechanism. Given reference $r(t)$,
measured output $y(t)$, and tracking error $e(t)=r(t)-y(t)$, proportional action
responds to present error, integral action addresses persistent accumulated
error, and derivative action anticipates rapid error change
\cite{astrom2001future}. These roles are meaningful in a physical loop with a
plant, time scale, actuation, and stability analysis.

Mapping PID structure to classifier fusion requires narrower language. Two
detector outputs are measurements, not a reference and controlled state.
Changing their weight does not alter the program being observed or close a
physical loop. In our formulation, static--dynamic disagreement is a surrogate
error and the controller output is a bounded mixture weight. The P/I/D terms
then provide hypotheses about useful memory: P reacts to current disagreement,
I represents persistent disagreement across prefixes, and D represents a change
in disagreement. No stability guarantee follows from this analogy. Validation
and component ablation decide whether any term improves prediction.

This framing addresses a common methodological gap. A controller name can make a
fusion rule appear inherently adaptive, even when each sample provides one
observation or the derivative rarely changes a decision. We retain P/PI/PID as
nested, testable models and report a negative I/D result. The contribution is
therefore an auditable translation and empirical simplification of a
control-inspired rule, rather than a claim that classical control theory proves
malware-classification performance.

\subsection{Active Learning and Batch Acquisition}

Pool-based active learning alternates model fitting, query selection, oracle
annotation, and retraining. Uncertainty sampling is widely used because it is
simple and model compatible \cite{cohn1996active,settles2009active}. In batch
settings, however, the highest-scoring samples can be redundant. Core-set
acquisition emphasizes representation coverage \cite{sener2018active}; ensemble
methods use predictive disagreement \cite{beluch2018power}; BatchBALD uses joint
mutual information \cite{kirsch2019batchbald}; and BADGE clusters uncertain
gradient embeddings \cite{ash2020deep}.

These methods are not equally portable to boosted trees. BatchBALD and BADGE
were designed around Bayesian or differentiable deep models, although their
principles motivate tree-compatible diversity and committee baselines. Our
measured policies include random sampling, uncertainty, two explanation-signal
ablations, and a three-signal combination. They do not include uncertainty plus
k-center diversity or a calibrated committee. We consequently test whether the
specific TreeSHAP summaries add value to the retained baselines, without claiming
general dominance over active learning.

Explanations may enter active learning as annotator-facing context or as an
algorithmic ranking signal. Explainable active learning has studied how local
explanations affect annotator experience \cite{ghai2021explainable}. Our use is
algorithmic: signed TreeSHAP summaries are computed before annotation and enter
the acquisition score. A balanced or dispersed attribution vector is not equal
to expected information gain; it is a falsifiable heuristic evaluated through
matched learning curves, target attainment, low-FPR summaries, and computation.

\subsection{Attribution and Explanation Validation}

LIME estimates a local surrogate \cite{ribeiro2016why}; SHAP unifies additive
feature attribution through Shapley values \cite{lundberg2017unified}; and
TreeSHAP specializes the calculation to tree ensembles
\cite{lundberg2020local}. Explanations can still be insensitive to model
parameters, manipulated, or misread as causal mechanisms
\cite{adebayo2018sanity,slack2020fooling}. Security inputs add domain constraints
and adversarial considerations \cite{warnecke2020evaluating}.

Deletion tests ask whether changing highly ranked coordinates alters model
output. Their conclusion depends on the intervention: independent replacement
can move a vector off the data manifold, and a large effect can reflect
perturbation magnitude as well as ranking quality
\cite{hooker2019benchmark,yeh2019infidelity}. Training-derived mean and median
baselines avoid assuming that zero is a semantically neutral PE value, but they
do not guarantee an executable-valid binary. Equal-count random masks control
the number of replaced coordinates, though they may select features unused by a
tree ensemble. Rank stability offers a complementary view by testing sensitivity
to model seeds and data folds.

\begin{table*}[t]
\centering
\caption{Positioning of the three components. “Open comparison” identifies an
evidence gate rather than a completed experiment.}
\label{tab:related}
\scriptsize
\setlength{\tabcolsep}{2.5pt}
\begin{tabular}{p{2.8cm}p{3.4cm}p{2.8cm}p{2.8cm}p{3.2cm}}
\toprule
Area & Representative approach & Strength & Limitation relevant here & Role in this study \\
\midrule
Static detection & EMBER/boosted trees \cite{anderson2018ember,ke2017lightgbm} & Early, scalable PE scoring & Static distribution shift and manipulation & AL/fidelity benchmark and static fusion channel \\
Dynamic detection & Sandbox/API behavior \cite{sgandurra2016automated} & Observes executed behavior & Environment and observation dependence & Event-indexed fusion channel \\
Fusion & Average/fixed/learned stacking & Transparent or data-adaptive combination & Fixed rules ignore disagreement; stacking can hide sequential behavior & Measured baselines for P/PI/PID \\
PID control & P/I/D feedback \cite{astrom2001future} & Interpretable current, accumulated, and rate terms & Classifier fusion has no physical plant or stability guarantee & Control-inspired nested ablation \\
Active learning & Uncertainty \cite{cohn1996active} & Simple and strong & Margin omits evidence structure & Measured baseline \\
Batch diversity & Core-set/committee \cite{sener2018active,beluch2018power} & Reduces redundancy or measures disagreement & Absent from retained artifact & Open comparison \\
Explanation guidance & Explainable AL \cite{ghai2021explainable} & Adds model rationale to selection/annotation & Explanation signal need not equal information gain & TreeSHAP summaries enter query score \\
XAI evaluation & Sanity/deletion tests \cite{adebayo2018sanity,hooker2019benchmark} & Tests model sensitivity & Intervention may be off-manifold & Mean/median deletion plus stability \\
\bottomrule
\end{tabular}
\end{table*}

Our novelty claim is correspondingly specific. We do not claim that static--
dynamic fusion, PID structure, active learning, or SHAP is individually new. We
contribute a scoped evidence-allocation formulation, a paired P/PI/PID fusion
ablation that supports simplification, a signed-mass TreeSHAP acquisition family,
and quantitative validation of the attribution ranking. The component-wise
protocol makes clear which links are empirically measured and which remain an
integration hypothesis.
